% GENERATED FILE. Edit canonical lessons, then run npm run generate:books.
% canonical-source-hash: fnv1a64:4b5e89ccdcb61528
% canonical-lessons: IT-C363-cittadino, IT-C363-cittadinanza, IT-C363-diritto, IT-C363-liberta, IT-C363-uguaglianza

\chapter{Cittadini}
\label{ch:it-cittadini}
\hlchaptermodality{363}

\begin{chapteropening}
\textbf{By the end of this chapter:} \emph{I can talk about citizens, citizenship, rights, freedom and equality.}

Complete the last lesson of chapter 363: I can talk about citizens, citizenship, rights, freedom and equality.
\end{chapteropening}

\section[il cittadino]{il cittadino — the citizen}

\label{lesson:IT-C363-cittadino}



\begin{quote}
Before the new one: say the Italian for efficient, then the Italian for the productivity.
\end{quote}

\subsection*{You'll want to know: il cittadino}
\textbf{il cittadino} — \textquotedblleft{}the citizen\textquotedblright{}.

\begin{quote}
Ogni cittadino può votare a diciotto anni.
\end{quote}

\begin{quote}
\emph{Every citizen can vote at eighteen.}
\end{quote}

As an adjective it means of the city: \emph{il traffico cittadino}.

\begin{cousinweb}
\emph{cittadino} is built on \emph{città}, older \emph{cittade}, inherited from Latin \emph{cīvitās, cīvitātis} \textquotedblleft{}citizenship, the body of citizens, a city\textquotedblright{} (EtymDB).
\end{cousinweb}

\begin{grammarlens}[title={What you've built}]
One more word for telling, explaining and describing at B1.
\end{grammarlens}

\subsection*{Your turn}
\begin{itemize}
\raggedright
  \item \emph{Say it:} \emph{il cittadino}
  \item \emph{Say it:} \emph{il cittadino}, once more
  \item \emph{Say it:} a sentence of your own with \emph{il cittadino}
  \item \emph{Recall:} say the Italian for efficient, then the Italian for the productivity, then say \emph{il cittadino} again
\end{itemize}

\begin{tcolorbox}[breakable,colback=teal!4,colframe=teal!35!black,title={Before you move on}]
What does il cittadino mean? (\textquotedblleft{}The citizen\textquotedblright{}.) Say it once more.
\end{tcolorbox}
\section[la cittadinanza]{la cittadinanza — the citizenship}

\label{lesson:IT-C363-cittadinanza}



\begin{quote}
Before the new one: say the Italian for the productivity, then the Italian for the citizen.
\end{quote}

\subsection*{You'll want to know: la cittadinanza}
\textbf{la cittadinanza} — \textquotedblleft{}the citizenship\textquotedblright{}.

\begin{quote}
Dopo dieci anni ha chiesto la cittadinanza.
\end{quote}

\begin{quote}
\emph{After ten years she applied for citizenship.}
\end{quote}

It can also mean the citizens of a town as a group.

\begin{cousinweb}
\emph{cittadinanza} is built on \emph{cittadino} \textquotedblleft{}citizen\textquotedblright{}, from \emph{città}, which is inherited from Latin \emph{cīvitās} \textquotedblleft{}citizenship, city\textquotedblright{} (EtymDB).
\end{cousinweb}

\begin{grammarlens}[title={What you've built}]
One more word for telling, explaining and describing at B1.
\end{grammarlens}

\subsection*{Your turn}
\begin{itemize}
\raggedright
  \item \emph{Say it:} \emph{la cittadinanza}
  \item \emph{Say it:} \emph{la cittadinanza}, once more
  \item \emph{Say it:} a sentence of your own with \emph{la cittadinanza}
  \item \emph{Recall:} say the Italian for the productivity, then the Italian for the citizen, then say \emph{la cittadinanza} again
\end{itemize}

\begin{tcolorbox}[breakable,colback=teal!4,colframe=teal!35!black,title={Before you move on}]
What does la cittadinanza mean? (\textquotedblleft{}The citizenship\textquotedblright{}.) Say it once more.
\end{tcolorbox}
\section[il diritto]{il diritto — the right; law (as a subject)}

\label{lesson:IT-C363-diritto}



\begin{quote}
Before the new one: say the Italian for the citizen, then the Italian for the citizenship.
\end{quote}

\subsection*{You'll want to know: il diritto}
\textbf{il diritto} — \textquotedblleft{}the right; law (as a subject)\textquotedblright{}.

\begin{quote}
Tutti hanno diritto a un avvocato.
\end{quote}

\begin{quote}
\emph{Everyone has the right to a lawyer.}
\end{quote}

\emph{Avere diritto a} — to be entitled to. \emph{Studiare diritto} — to study law. As an adjective it means straight.

\begin{cousinweb}
\emph{diritto} is inherited from Vulgar Latin \emph{dērēctus}, a form of Latin \emph{dīrēctus} \textquotedblleft{}straight\textquotedblright{} (EtymDB), from \emph{dīrigere} \textquotedblleft{}to set straight\textquotedblright{}, built on \emph{regere} \textquotedblleft{}to rule\textquotedblright{}: what is straight is what is right.
\end{cousinweb}

\begin{grammarlens}[title={What you've built}]
One more word for telling, explaining and describing at B1.
\end{grammarlens}

\subsection*{Your turn}
\begin{itemize}
\raggedright
  \item \emph{Say it:} \emph{il diritto}
  \item \emph{Say it:} \emph{il diritto}, once more
  \item \emph{Say it:} a sentence of your own with \emph{il diritto}
  \item \emph{Recall:} say the Italian for the citizen, then the Italian for the citizenship, then say \emph{il diritto} again
\end{itemize}

\begin{tcolorbox}[breakable,colback=teal!4,colframe=teal!35!black,title={Before you move on}]
What does il diritto mean? (\textquotedblleft{}The right; law (as a subject)\textquotedblright{}.) Say it once more.
\end{tcolorbox}
\section[la libertà]{la libertà — the freedom, the liberty}

\label{lesson:IT-C363-liberta}



\begin{quote}
Before the new one: say the Italian for the citizenship, then the Italian for the right.
\end{quote}

\subsection*{You'll want to know: la libertà}
\textbf{la libertà} — \textquotedblleft{}the freedom, the liberty\textquotedblright{}.

\begin{quote}
La libertà di parola è importante.
\end{quote}

\begin{quote}
\emph{Freedom of speech is important.}
\end{quote}

\emph{In libertà} — free, at large. Stress on the last syllable.

\begin{cousinweb}
\emph{libertà} is a learned borrowing from Latin \emph{lībertās, lībertātis} \textquotedblleft{}freedom\textquotedblright{} (Lewis \& Short), from \emph{līber} \textquotedblleft{}free\textquotedblright{}, a different word from \emph{liber} \textquotedblleft{}book\textquotedblright{}.
\end{cousinweb}

\begin{grammarlens}[title={What you've built}]
One more word for telling, explaining and describing at B1.
\end{grammarlens}

\subsection*{Your turn}
\begin{itemize}
\raggedright
  \item \emph{Say it:} \emph{la libertà}
  \item \emph{Say it:} \emph{la libertà}, once more
  \item \emph{Say it:} a sentence of your own with \emph{la libertà}
  \item \emph{Recall:} say the Italian for the citizenship, then the Italian for the right, then say \emph{la libertà} again
\end{itemize}

\begin{tcolorbox}[breakable,colback=teal!4,colframe=teal!35!black,title={Before you move on}]
What does la libertà mean? (\textquotedblleft{}The freedom, the liberty\textquotedblright{}.) Say it once more.
\end{tcolorbox}
\section[l'uguaglianza]{l'uguaglianza — the equality}

\label{lesson:IT-C363-uguaglianza}



\begin{quote}
Before the new one: say the Italian for the right, then the Italian for the freedom.
\end{quote}

\subsection*{You'll want to know: l'uguaglianza}
\textbf{l'uguaglianza} — \textquotedblleft{}the equality\textquotedblright{}.

\begin{quote}
Lottiamo per l'uguaglianza tra uomini e donne.
\end{quote}

\begin{quote}
\emph{We fight for equality between men and women.}
\end{quote}

Also spelled \emph{eguaglianza}.

\begin{cousinweb}
\emph{uguaglianza} is built on \emph{uguale} \textquotedblleft{}equal\textquotedblright{}, inherited from Latin \emph{aequālis} \textquotedblleft{}equal\textquotedblright{}, from \emph{aequus} \textquotedblleft{}level, even\textquotedblright{} (Lewis \& Short).
\end{cousinweb}

\begin{grammarlens}[title={What you've built}]
One more word for telling, explaining and describing at B1.
\end{grammarlens}

\subsection*{Your turn}
\begin{itemize}
\raggedright
  \item \emph{Say it:} \emph{l'uguaglianza}
  \item \emph{Say it:} \emph{l'uguaglianza}, once more
  \item \emph{Say it:} a sentence of your own with \emph{l'uguaglianza}
  \item \emph{Recall:} say the Italian for the right, then the Italian for the freedom, then say \emph{l'uguaglianza} again
\end{itemize}

\begin{tcolorbox}[breakable,colback=teal!4,colframe=teal!35!black,title={Before you move on}]
What does l'uguaglianza mean? (\textquotedblleft{}The equality\textquotedblright{}.) Say it once more.
\end{tcolorbox}
